\section*{Capítulo \ref{ch:b3:x-ray-diffraction} --- Cristalografía y difracción de rayos X}

\begin{solution}{exo:b3:x-ray-diffraction:1}
$d = a/\sqrt N$: $d_{111} = \qty{325.66}{pm}$, $d_{200} = \qty{282.03}{pm}$, $d_{220} =
\qty{199.42}{pm}$.
\end{solution}

\begin{solution}{exo:b3:x-ray-diffraction:2}
FCC: las primeras reflexiones son (111) y (200). $\sin\theta = \lambda\sqrt N/2a$: $2\theta =
\ang{43.32}$ y \ang{50.45}.
\end{solution}

\begin{solution}{exo:b3:x-ray-diffraction:3}
P: las siete. I ($h + k + l$ par): 110, 200, 211, 220. F (no mixtos): 111, 200, 220.
\end{solution}

\begin{solution}{exo:b3:x-ray-diffraction:4}
$\mathbf a^* = \mathbf e_x/a$, $\mathbf b^* = \mathbf e_y/b$: una red rectangular de lados $1/a$
y $1/b$; el lado largo de la red directa pasa a ser el corto.
\end{solution}

\begin{solution}{exo:b3:x-ray-diffraction:5}
Razones de $\sin^2\theta$: $1 : 2 : 3 : 4 : 5$, es decir, $N = 2, 4, 6, 8, 10$: centrada en el cuerpo (110,
200, 211, 220, 310). A partir de la primera raya, $d_{110} = \lambda/2\sin(\ang{20.135}) =
\qty{223.77}{pm}$ y $a = d\sqrt2 = \qty{316.5}{pm}$: el wolframio.
\end{solution}

\begin{solution}{exo:b3:x-ray-diffraction:6}
$Z = \rho N_Aa^3/M = 1.99 \times 6.022\times10^{23} \times (6.2929\times10^{-8})^3/74.6 = 4.0$.
\end{solution}

\begin{solution}{exo:b3:x-ray-diffraction:7}
$\beta = \qty{0.40}{\degree} = \qty{6.98e-3}{rad}$, $\cos\theta = \cos\ang{15.85} = 0.962$:
$L = 0.9 \times 0.15406/(6.98\times10^{-3} \times 0.962) = \qty{21}{nm}$.
\end{solution}

\begin{solution}{exo:b3:x-ray-diffraction:8}
$F_{hkl} = f_{\ce{Cs+}} + f_{\ce{Cl-}}(-1)^{h+k+l}$: $F_{100} = 54 - 18 = 36$, $F_{110} = 72$:
intensidades en la razón 1 : 4. El centro está ocupado por un ion distinto del de
los vértices: la red es primitiva (una red centrada en el cuerpo exige contenidos
idénticos en ambos puntos), y 100 está presente.
\end{solution}

\begin{solution}{exo:b3:x-ray-diffraction:9}
Cada átomo tiene una copia desplazada en $(\frac12,\frac12,0)$, lo que multiplica $F$ por $1 +
\eu^{\iu\pi(h+k)} = 1 + (-1)^{h+k}$, nulo para $h + k$ impar.
\end{solution}

\begin{solution}{exo:b3:x-ray-diffraction:10}
Los cuasicristales tienen orden de largo alcance sin periodicidad: no tienen
red, de modo que la restricción, que concierne a las redes, no se les aplica.
Sus manchas de difracción nítidas llevaron a definir un cristal como todo sólido cuyo
diagrama de difracción es esencialmente discreto, periódico o no.
\end{solution}

\begin{solution}{exo:b3:x-ray-diffraction:11}
El deslizamiento $c$ transforma $(x, y, z)$ en $(x, \frac12 - y, \frac12 + z)$. Para $k = 0$, los dos
factores de fase son $\eu^{2\pi\iu(hx + lz)}$ y $\eu^{2\pi\iu(hx + lz + l/2)} = (-1)^l
\eu^{2\pi\iu(hx+lz)}$: su suma se anula para $l$ impar.
\end{solution}

\begin{solution}{exo:b3:x-ray-diffraction:12}
Una inversión, un \omterm{def:b3:point-groups:mirror}{plano de simetría} o un deslizamiento transforman una molécula en su imagen especular, de modo que un
cristal que contuviera uno contendría ambos enantiómeros. Un enantiómero puro
solo cristaliza en los 65 \omterm{def:b3:x-ray-diffraction:space-group}{grupos espaciales} construidos con rotaciones, \omterm{def:b3:x-ray-diffraction:space-group}{ejes helicoidales}
y traslaciones (grupos de Sohncke), como \textit{P}$2_12_12_1$. Por la ley de Friedel, el
diagrama de un enantiómero es igual al del otro; la configuración
absoluta se obtiene de las pequeñas desviaciones causadas por la dispersión anómala
de los átomos más pesados.
\end{solution}

\begin{solution}{pb:b3:x-ray-diffraction:1}
\textbf{1.} $\theta = \ang{14.171}$, $d = \lambda/2\sin\theta = \qty{314.64}{pm}$.
\textbf{2.} 314.64, 222.49, 181.66, 157.32, 140.71, \qty{128.45}{pm}.
\textbf{3.} 1.000, 2.000, 3.000, 4.000, 5.000, 6.000.
\textbf{4.} 1, 2, 3, 4, 5, 6.
\textbf{5.} Sí, como las (100), (110), (111), (200), (210), (211) de una celda cúbica primitiva
con $a' = \qty{314.6}{pm}$.
\textbf{6.} $N = 7$ (y 15, 23\dots), que no es suma de tres cuadrados; la
primera diferencia aparecería en la octava raya. Seis rayas no permiten distinguir una celda P
de arista $a'$ de una celda F de arista $2a'$.
\textbf{7.} $N = 4, 8, 12, 16, 20, 24$: (200), (220), (222), (400), (420), (422).
\textbf{8.} $a = 2d_{200} = \qty{629.3}{pm}$.
\textbf{9.} \ce{KCl} (\qty{629.29}{pm}).
\textbf{10.} $F = 4[f(\ce{K+}) - f(\ce{Cl-})]$ para índices todos impares, y los dos iones tienen 18
electrones cada uno: las amplitudes se anulan.
\textbf{11.} Sí en ambos casos: \ce{Na+} (10) y \ce{Cl-} (18), \ce{K+} (18) y \ce{Br-} (36)
dispersan de forma distinta; la (111) del \ce{NaCl} estaría a \ang{27.36}, y la del \ce{KBr} a \ang{23.38},
ambas dentro del intervalo registrado.
\textbf{12.} Los rayos X ven \ce{K+} y \ce{Cl-} casi idénticos; unos iones idénticos en una
disposición de sal gema forman una red cúbica simple de arista $a/2$, una celda aparente.
\textbf{13.} Cuatro \ce{KCl}.
\textbf{14.} $\rho = 4 \times 74.6/(6.022\times10^{23} \times (6.293\times10^{-8})^3) =
\qty{1.99}{g/cm^3}$.
\textbf{15.} Seis y seis (octaedros del otro ion).
\textbf{16.} $a/2 = \qty{314.6}{pm}$.
\textbf{17.} La (440), $N = 32$, a \ang{87.65} (pues faltan las (333)/(511), de índices todos impares).
\textbf{18.} Las intensidades cambian con la orientación preferente de los cristalitos,
la absorción, la disminución de los factores de dispersión y la agitación térmica; las posiciones
solo dependen de la red.
\textbf{19.} A partir de $d = \lambda/2\sin\theta$, $\Delta d/d = -\cot\theta\,\Delta\theta$: el mismo error angular
da un error relativo menor a $\theta$ grande.
\textbf{20.} $d_{420} = \lambda/2\sin\ang{33.190} = \qty{140.71}{pm}$, $a = d\sqrt{20} =
\qty{629.3}{pm}$.
\textbf{21.} $L = 0.9 \times 0.15406/(4.36\times10^{-3} \times \cos\ang{33.19}) = \qty{38}{nm}$.
\textbf{22.} No: $\beta = \ang{0.0095}$, muy por debajo de una anchura instrumental típica.
\textbf{23.} \textbf{$a = \qty{629.3}{pm}$: el polvo es cloruro de potasio}, cuyas reflexiones
impares silencian dos iones con el mismo número de electrones.
\end{solution}
